\documentclass{ltjbook}
\usepackage{docmute} 
\usepackage{../myStyle} 

\begin{document} 
\chapter{後期のまとめ}
\label{m30_SecondSemester} 
後期のテーマを振り返ってみます。

前期は記述統計，線形モデル(要因計画)，誤差としての確率が導入されました。
後期は\keytermL{確率}{かくりつ}が中心的なテーマになっています。というのも\keytermL{推測統計学}{すいそくとうけいがく}の領域に入ったからですね。
前期の要因計画で，群平均の差を見ることで，介入の効果などが表現できることが論じられました。科学的実験はもちろん，カウンセリングなどの臨床的介入においても，その効果があったのかなかったのかを評価することは重要です。それが目の前の被験者，クライアントだけでなく，どれほど一般化できるかということが重要なのです。そのためには，目の前のデータに対して\keytermL{母集団}{ぼしゅうだん}と\keytermL{標本}{ひょうほん}という区別をつけ，標本から母集団を推測する必要があるわけです。

標本から母集団を考えてものをいうためには，\keytermL{推定}{すいてい}と\keytermL{検定}{けんてい}という考え方を使います。

推定は，標本から母集団の状態を記述することを考えるものです。母平均，母相関，母比率など\keytermL{母数}{ぼすう}がどうなっているのかについて，手元の標本平均，標本相関係数，標本分布などからその数字を見積もるわけですね。方法としては，モーメント法，最尤法，ベイズ法の3つが考えられているのですが，モーメント法を使ったものが最も広く利用されています。

\keytermL{モーメント法}{もーめんとほう}の「もーめんと」とは，積率という訳語もあるのですが，要するに平均，分散，歪度，尖度などのことです。標本平均や標本分散から，母平均や母分散を推定するというだけでも，少数のデータから全体像を推測するというなかなか難しいチャレンジをしているんだ，ということをまずは踏まえて欲しいところです。推測統計学というのはことほど左様に，そもそも問題を解決するためのヒントが少なすぎるという難易度の高いものでもあるのです。

標本分散は母分散に一致せず，ずれてしまいますので\keytermL{不偏分散}{ふへんぶんさん}という計算方法をしなければなりませんでした。標本平均がどの程度の散らばりをするかは，標本分の標準偏差，すなわち\keytermL{標準誤差}{ひょうじゅんごさ}を使って表現できますが，母SDがわからないのであればそのままでは使えない。ということで不偏SDとt分布をつかって推定する必要があるのでした。\keytermL{区間推定}{くかんすいてい}はベイズ推定のそれとは確率の考え方が違っているので注意が必要ですが，要因計画で介入の効果を考えるのが，群平均に基づくものだったわけですから，これを使って効果の大きさを見積もるのだ，という方針をしっかり思い出してください。

もう1つの推測統計学の用法，統計的な\textbf{検定}というのは，こうした推定と判断がセットになった考え方です。差があるのかないのか，効果があるのかないのか，という判断をする必要があったときに用いられます。手続きは機械的に進められるカッチリしたもので，\keytermL{帰無仮説}{きむかせつ}と\keytermL{対立仮説}{たいりつかせつ}という2つの仮説による対決で勝敗を決するというものです。判断基準である\keytermL{有意水準}{ゆういすいじゅん}や，\keytermL{検定力}{けんていりょく}，確率的判断の正しさなど細々したルールがあるので，しっかりとその手続きについても理解しておいてください。また，相関係数の判断についてはt分布，要因計画の判断についてはF分布など，対応する統計量が何になるのかとか，それぞれの帰無仮説はどのように設定するのか，効果量はどのように見積もるのかといったことも，実践上は重要なポイントになってきます。

検定は誤用や悪用もあるのですが，そもそも母数を推定したい，科学的に判断したいというニーズに基づいていたことをしっかり思い出せば，こうした問題は幾らか軽減されるかもしれません。推定が目的であることを思い出せば，差があったのかなかったのか，モデルの勝敗はどうだったのかという一点にこだわりすぎることもなくなるでしょう。また「有意になりさえすればよい」「有意な検定結果が良い論文」という価値判断に陥ることもなくなるでしょう。科学論文は正しいかどうかが重要で，嘘をついてまで有意差を見出すことが目的ではないからです。こうした考え方は，モーメント法による有意差検定だけが母集団の推定方法だと思っていることが原因かもしれません。

最尤法やベイズ法では，確率モデルを元に母集団を推測する方法です。確率関数を直接用いることになるので，数式や計算方法がややこしいように思えてきますが，そもそも確率は「個々の事例はわからないけれども全体的な傾向はわかる」という\underline{偶然をハンドリングする}ための秘策でもあるのです。理論的にはこうなるはず，でも完全に制御できない要素がある，それがどの程度かを表現する，そういう方法が確率です。個々のデータは確率分布から生み出されているんだと考えると，得られたデータがどういう確率分布から出てきているんだろうか，ということが問題になってきます。得られたデータに合うように確率分布のパラメータを推定する，それが\keytermL{最尤法}{さいゆうほう}のやりかたなのでした。

尤度はデータが得られる確率から計算するもので，尤度関数自体は確率分布のように解釈することはできません。形は似ているのですが。
これを確率のように解釈するには，ベイズのやり方を使う必要があります。注意すべきは，ベイズ法では確からしさ，確信の強さを確率分布で表現するという解釈に立っているということです。確率分布は信念の強さ，確率は\keytermL{主観確率}{しゅかんかくりつ}というべきものとして理解した方がいいでしょう。であれば，母数がどのあたりにあると思われるかを表現するのがベイズ推定の結果であり，結果を表す\keytermL{事後分布}{じごぶんぷ}はその名の通り，分布として区間の情報を持っているものになります。区間を使って推定する，判断するという枠組みは同じですが，帰無仮説検定のような架空の確率空間におけるp値を見るのではなく，母数の分布そのものを見ることができるので，より直感的な理解が可能です。実践上はJASPやStanなど特殊な統計環境を必要としますから，まだ一般に広まった考え方とは言えませんが，今後ますます活用されてくることになるでしょう。


心理学における統計の使い方は，「心」はもちろん「実験・介入の効果」といった目に見えないものを，なんとか測定し，推定し，心理学的考察の補助とするためのものです。
これから皆さんは，様々な専門的心理学の道に進んでいくことになりますが，そのどこにおいても方法論的基礎は必ず利用します。心理学を学べば学ぶほど，人間の直感的な理解や判断は間違えていることが多く，数学的あるいは理論的に正しい判断をするための補助的なツールが必要だということが身に染みてくると思います。そのための方法論ですから，決して蔑ろにすることなく，正しく身につけてもらいたいと思います。

どうしても数学的な表現をすることがありますので，その表面的な難しさに嫌気が差して目を逸らしてしまいたくなる気持ちもわかります。幸い，わからなければ立ち返って，わからなくなったところからやり直すことができるのが，数学のよいところです。私自身，間違えた理解や判断をしたことがないはずはなく，その度に学び直すということを繰り返しています。怖がらなくて結構ですから，「何のために統計分析をしているのか」を忘れずにこれからも学習し続けていってください。


データ解析応用の授業では，さらに進んだ理解，表現ができるようなことを学びます。基礎の授業では要因計画や帰無仮説検定など，決まり切った方法をなぞることを求めてきましたが，これから先は統計を駆使して，よりクリエイティブに表現することを考えます。応用まで履修していただくと，心理測定の理論や心を表現するための新しいモデル，心のメカニズムを数式やプログラミングで表現できるようになります。これからの社会はデータサイエンティストが求められていますが，人とデータ，データとその表現について心理統計は一日の長がありますので，ぜひ履修してみてください。もちろんデータサイエンティストが流行らなくなった時代になったとしても，人間とは何か，心とは何かという問いについて，その分析方法がもつ思想的背景からの観点を身につけておけば，人生をより豊かなものとして楽しむことができるでしょう。

\end{document} 
